\documentclass{situationist-preprint}

\title{The First Horn of Hempel's Dilemma}
\msid{MND-02-b}
\msfamily{MND-02}
\msdiscipline{Philosophy of Mind}
\mstype{Companion manuscript}
\mskeywords{physicalism; Hempel's dilemma; current physics; foreseeable extension; scientific revolutions}

\begin{document}
\maketitle

\begin{abstract}
The first horn of Hempel's dilemma says that physicalism defined by current physics is false, because current physics is incomplete and in part mistaken. MND-02 escapes it by defining the physical through current physics together with its foreseeable extensions. We examine the horn more closely. We show that it rests on an assumption, that the definition of \enquote{physical} must take over the content of current physics wholesale, which the definition of MND-02 rejects; that the anchoring in current physics needed to avoid the second horn survives the revision of current theories, including revolutionary revision; and that the current-physics response, which accepts the horn, concedes more than it needs to.
\end{abstract}

\section{The horn and its assumption}

\citet{Hempel1980} put the first horn briefly: physicalism formulated by reference to the physics of today is almost certainly false, since that physics is known to be incomplete and will be revised. The inference has a hidden premise.

\begin{principle}[Wholesale inheritance]\label{p:whole}
If \enquote{physical} is defined by reference to current physics, then a property is physical only if current physics posits it, and the laws and ontology of current physics are true.
\end{principle}

Given Principle~\ref{p:whole}, the falsity of current physics in any respect makes physicalism false or makes the physical empty. But the principle is not forced by the project of defining the physical through current physics. A definition can use a theory as a point of reference without asserting it. The definition of MND-02 does exactly this: a foreseeable extension must recover the predictions of current physics where current physics is empirically successful, and nothing more is inherited (MND-02-a, Definition~2.2(a)).

\begin{proposition}\label{p:escape}
Physicalism in the sense of MND-02 is not refuted by the incompleteness or partial falsity of current physics.
\end{proposition}

\begin{proof}
By Definition~2.3 of MND-02-a, a property is physical if it is a fundamental property of some foreseeable extension of current physics or wholly grounded in such properties. Definition~2.2 imposes on a foreseeable extension correspondence with the successful predictions of current physics, nomic integration and the absence of fundamental mentality. None of these conditions requires that current physics be complete or true. An extension that adds properties current physics lacks, or replaces its laws by others that agree with them within the domain of their success, satisfies all three. So the incompleteness or falsity of current physics does not entail that any property fails to be physical, and does not entail that physicalism is false.
\end{proof}

\section{Anchoring without inheritance}

The danger is that Proposition~\ref{p:escape} is bought by moving to the second horn: once current physics is not asserted, \enquote{physical} may seem to mean whatever a successor theory will posit. What prevents this is that two of the three conditions of Definition~2.2 of MND-02-a constrain content now. Correspondence fixes the domain in which any successor must agree with what we have; the exclusion of fundamental mentality fixes a kind of property no successor may posit as fundamental. The second of these is what makes physicalism non-trivial (MND-02-a, Proposition~3.5).

\begin{objection}
Scientific revolutions do not preserve the content of the theories they replace. The ontology of classical mechanics is not that of quantum theory. If the next revolution is as deep, correspondence to current physics anchors nothing.
\end{objection}

\begin{reply}
Correspondence as defined concerns predictions within domains of empirical success, not ontology. Revolutions in physics have preserved predictions of this kind: quantum mechanics recovers classical mechanics for large actions, and relativity recovers Newtonian gravitation for weak fields and low velocities. That ontology changes is consistent with the definition, which lets the ontology of the extension differ from ours. The objection would succeed only against a theory-based conception that inherited ontology, which is Principle~\ref{p:whole} again.
\end{reply}

\begin{objection}
The current-physics response accepts the horn and treats physicalism as a hypothesis known to be false in detail but held as the best approximation \citep{Melnyk1997}. This is more honest than a definition designed to avoid falsification.
\end{objection}

\begin{reply}
The current-physics response gives up the claim that physicalism is true and replaces it with the claim that it is a good approximation. That is a concession the dilemma does not compel, since the dilemma depends on Principle~\ref{p:whole}. The definition of MND-02 is not designed to avoid falsification: Proposition~3.5 of MND-02-a identifies the worlds at which physicalism is false, namely those with ungrounded mental properties. It avoids only the falsification that comes from the revision of physics, which a thesis about the mental should not be hostage to. As \citet{Wilson2006} argues, what physicalists care about is the exclusion of fundamental mentality, and that is what the definition secures.
\end{reply}

\section{Conclusion}

The first horn depends on the principle that a definition of the physical through current physics inherits its content wholesale. The definition of MND-02 uses current physics as an anchor for correspondence and inherits nothing else, so the horn does not apply, and the anchoring is strong enough, together with the exclusion of fundamental mentality, to keep the definition off the second horn. Physicalism so defined is true.

\begin{thebibliography}{9}\small
\bibitem[Hempel(1980)]{Hempel1980} Hempel, C.~G. (1980). Comments on Goodman's \emph{Ways of Worldmaking}. \emph{Synthese}, 45(2), 193--199.
\bibitem[Melnyk(1997)]{Melnyk1997} Melnyk, A. (1997). How to keep the \enquote{physical} in physicalism. \emph{The Journal of Philosophy}, 94(12), 622--637.
\bibitem[Wilson(2006)]{Wilson2006} Wilson, J. (2006). On characterizing the physical. \emph{Philosophical Studies}, 131(1), 61--99.
\end{thebibliography}

\end{document}
